\section{An almost Euclidean critical Sobolev inequality}
\label{sec:sobolev}

Throughout this section, $n>2$, and the sharp filling theorem in dimension
$n-1$ is assumed.  We work with the extremizing cycle $A\in\I_n(Y)$ of
Lemma~\ref{lem:extremizer}, normalized as in
\eqref{eq:normalization}.  In particular, $Y$ is compact and CAT(0),
$\mu=\|A\|$, $m=\M(A)>0$, and
\[
 q_0=\frac{n+1}{n},\qquad cq_0m^{1/n}=\frac1n.
\]
The covectors $Du$ and their norms are those of
Proposition~\ref{prop:chart-mass}.  Set
\begin{equation}
 \beta=\frac1{n-2},\qquad p=\frac{2n}{n-2},\qquad
 S_*=ns_n^{2/n},\qquad
 E(u)=\int_Y|Du|^2\dd\mu.
 \label{eq:sobolev-notation}
\end{equation}
Unless another measure is specified, function norms in this section are
taken with respect to $\mu$.

Our objective is an almost sharp critical estimate
\[
 (S_*-\epsilon)\|u\|_p^2
 \le4\beta E(u)+C_\epsilon\|u\|_2^2
\]
uniformly over Lipschitz $u$. We first use the dimension $n-1$ filling
theorem to control the boundary mass of small restrictions of $A$.
Coarea and rearrangement transfer this estimate to functions with small
support; a level truncation then gives the displayed bound for every $u$.

\subsection{Small sets and intrinsic coarea}

\begin{proposition}[Uniform small-set isoperimetry]
\label{prop:small-set}
For every $\eta\in(0,1)$ there is a number $\delta\in(0,m/2)$, independent
of the Lipschitz function $u\colon Y\to\R$, with the following property.
Put
\[
 a(t)=\mu\{u>t\},\qquad
 U_t=A\restr\{u>t\},\qquad P(t)=\M(\partial U_t),
\]
where $U_t$ and its boundary are integral for almost every $t$.
At almost every level for which $a(t)\le\delta$,
\begin{equation}
 P(t)\ge(1-\eta)n\omega_n^{1/n}a(t)^{(n-1)/n}.
 \label{eq:small-set}
\end{equation}
\end{proposition}

\begin{proof}
Fix a level at which $U_t$ is integral, and write $a=a(t)$ and $P=P(t)$.
Its boundary is an integral $(n-1)$-cycle.  By the induction hypothesis
there is an integral $n$-current $R$ in $Y$ satisfying
\[
 \partial R=\partial U_t,\qquad
 d:=\M(R)\le
 \left(\frac{P}{n\omega_n^{1/n}}\right)^{n/(n-1)}.
\]
Indeed, the coefficient on the right is the sharp filling coefficient
in dimension $n-1$.  The currents $B=A-U_t+R$ and $A-B=U_t-R$ are
integral cycles.  Restriction decomposes the mass measure exactly, so
\[
 \M(B)\le m-a+d,\qquad \M(A-B)\le a+d.
\]
No disjointness involving $R$ is needed for these upper bounds.

Suppose first that $d<a\le m/2$.  The extremizer inequality
\eqref{eq:extremizer} and the coarse filling estimate in
Theorem~\ref{thm:current-background} give
\begin{align*}
 cq_0(m/2)^{1/n}(a-d)
 &\le c\bigl(m^{q_0}-(m-a+d)^{q_0}\bigr)\\
 &\le V(A-B)
 \le K_n(a+d)^{q_0}
 \le K_n(2a)^{q_0}.
\end{align*}
The first inequality follows by integrating the derivative of
$s^{q_0}$ over $[m-a+d,m]\subset[m/2,m]$.  Thus
\begin{equation}
 d\ge a-Ca^{1+1/n},\qquad
 C=\frac{K_n2^{q_0}}{cq_0(m/2)^{1/n}}
   =nK_n2^{1+2/n}.
 \label{eq:small-set-replacement}
\end{equation}
If $d\ge a$, the same lower bound holds automatically.  Choose
$0<\delta<m/2$ so small that
\[
 C\delta^{1/n}\le1-(1-\eta)^{n/(n-1)}.
\]
For $0<a\le\delta$, combining \eqref{eq:small-set-replacement} with the
upper bound for $d$ yields
\[
 P\ge n\omega_n^{1/n}a^{(n-1)/n}
           (1-Ca^{1/n})^{(n-1)/n}
   \ge (1-\eta)n\omega_n^{1/n}a^{(n-1)/n}.
\]
The case $a=0$ is immediate.  All constants used to choose $\delta$
depend only on the fixed extremizer, the dimension, and $\eta$.
\end{proof}

We now have the Euclidean perimeter coefficient, up to an arbitrarily
small loss, for every sufficiently small superlevel restriction. To turn
this into an energy inequality, coarea must retain the chart derivative
$|Du|$. Its uniform upper bound $\Lip(u)$ would lose the required
Dirichlet energy. The next lemma recovers each slice's total mass from
countably many scalar evaluations, so that the integrated cycle identity
can be used at almost every level.

\begin{lemma}[A countable family of mass tests]
\label{lem:countable-mass-tests}
Let $Z$ be a nonempty compact metric space and let $k\ge1$.  There is a
countable collection $\mathcal T_k(Z)$ of finite lists
\[
 \bigl((b_1,\pi^1),\ldots,(b_N,\pi^N)\bigr),
\]
where the $b_j$ are Lipschitz, $\sum_j|b_j|\le1$ pointwise, and each
$\pi^j$ is a $k$-tuple of $1$-Lipschitz functions, such that every
$k$-current $C$ of finite mass on $Z$ satisfies
\begin{equation}
 \M(C)=\sup_{\mathcal T\in\mathcal T_k(Z)}
                  \sum_{(b,\pi)\in\mathcal T}C(b,\pi).
 \label{eq:countable-mass-tests}
\end{equation}
The collection is independent of $C$.
\end{lemma}

\begin{proof}
Fix a point $z_0\in Z$ and normalize each differentiated test by
$\pi_i(z_0)=0$; adding constants does not change its current evaluation.
The normalized $1$-Lipschitz functions form a compact metric space in
the uniform topology, by the Arzel\`a--Ascoli theorem.  Choose a
countable dense family of normalized $k$-tuples
$(\pi^j)_{j\ge1}$.

Write $\lambda=\|C\|$.  For each tuple $\pi$, the first-slot extension
to bounded Borel functions gives a finite signed measure
$b\mapsto C(b,\pi)$ with density $\rho_\pi$ relative to $\lambda$,
where $|\rho_\pi|\le1$ almost everywhere.  Choose measurable versions
for the countable family and put
\[
 \rho(z)=\sup_{j\ge1}|\rho_{\pi^j}(z)|.
\]
This satisfies $0\le\rho\le1$.  If a normalized tuple $\pi$ is
approximated uniformly by tuples in the family, current continuity,
with their common Lipschitz bound, shows for every Lipschitz $b$ that
\[
 |C(b,\pi)|\le\int_Z|b|\rho\dd\lambda.
\]
Thus $\rho\lambda$ is a mass-controlling measure for all normalized
$1$-Lipschitz tuples and hence, by rescaling the differentiated slots,
for all tests.  Minimality of the mass measure gives
$\lambda\le\rho\lambda$.  It follows that $\rho=1$ almost everywhere.

For any finite set of the tuples, select at each point the first index
where the largest $|\rho_{\pi^j}|$ is attained and select its sign.
This produces bounded Borel functions $b_j$ with
$\sum_j|b_j|\le1$ and
\[
 \sum_j C(b_j,\pi^j)
   =\int_Z\max_j|\rho_{\pi^j}|\dd\lambda.
\]
These coefficients can be replaced, with an arbitrarily small error
in the displayed sum, by Lipschitz coefficients satisfying the same
constraint.  To see this, regularity of the finite Borel measure
$\lambda$ gives Lipschitz approximations to each $b_j$ in
$L^1(\lambda)$.  For example, an indicator is approximated between a
compact set and an open neighborhood by a Lipschitz function obtained
from distances to those sets, and simple functions suffice for the
general case.  Apply to the vector of approximations the Lipschitz map
\[
 (v_1,\ldots,v_N)\longmapsto
 \frac{(v_1,\ldots,v_N)}{\max\{1,\sum_j|v_j|\}}.
\]
This retains $L^1$ convergence and enforces the constraint.
Monotone convergence of the finite maxima to $\rho=1$ now proves that
the supremum over \emph{all} admissible finite Lipschitz lists equals
$\M(C)$; the reverse inequality follows directly from the mass bound.

It remains to make the lists countable without depending on $C$.
For each fixed list length, the admissible lists of coefficients and
normalized tuples form a separable metric space in the uniform
topology: they are a subspace of a finite product of $C(Z)$, which is
separable because $Z$ is compact metric.  Choose a countable dense
subset within that admissible space and take the union over list
lengths.  An admissible list can be approximated by these lists with
its constraint and differentiated Lipschitz bounds preserved.
The first-slot errors are controlled by their uniform norms times
$\M(C)$, and convergence in the differentiated slots follows from
current continuity.  Hence the countable collection has the same
supremum.  A common Lipschitz bound on the first coefficients is
neither asserted nor needed.
\end{proof}

\begin{lemma}[Intrinsic coarea inequalities]
\label{lem:intrinsic-coarea}
For a real Lipschitz function $u$ on $Y$, put
\[
 a(t)=\mu\{u>t\},\qquad
 e(t)=\int_{\{u>t\}}|Du|^2\dd\mu,
 \qquad P(t)=\M\bigl(\partial(A\restr\{u>t\})\bigr)
\]
at the almost every level where the restriction is integral.  Then
\begin{equation}
 P(t)^2\le(-a'(t))(-e'(t)),\qquad
 P(t)\le\Lip(u)(-a'(t))
 \quad\text{for almost every }t.
 \label{eq:intrinsic-coarea}
\end{equation}
The derivatives are the densities of the absolutely continuous parts
of the corresponding distribution measures; neither distribution
function is assumed absolutely continuous.
\end{lemma}

\begin{proof}
For a bounded open interval $I$, define the bounded Lipschitz function
\[
 H(s)=\int_I1_{\{t<s\}}\dd t.
\]
Let $b$ be Lipschitz and let $\pi$ be an $(n-1)$-tuple of
$1$-Lipschitz functions.  Restriction, the chart representation, and
Fubini give
\begin{equation}
 \int_I\partial U_t(b,\pi)\dd t
   =A(H(u),b,\pi)=-A(b,H(u),\pi).
 \label{eq:integrated-slice}
\end{equation}
For the first equality, the boundary evaluation has the measurable
representative $U_t(1,b,\pi)$, obtained by integrating the chart
determinants with the factor $1_{\{u>t\}}$.  It is integrable on bounded
intervals.  The second equality is the product rule applied to the
cycle identity $\partial A=0$.

The chain rule in the charts gives
\[
 D(H(u))=1_{\{u\in I\}}Du\qquad\mu\text{-almost everywhere}.
\]
There is no ambiguity at the two endpoints of $I$: the differential
of a Lipschitz scalar function vanishes almost everywhere on each
fixed level set.  This follows by applying Euclidean differentiation
at density points of that level set in each chart.  At points of the
charts the other differentiated covectors have norm at most one, so
the Euclidean determinant bound in
Proposition~\ref{prop:chart-mass} applies.

In particular, for an admissible finite list $\mathcal T$ as in
Lemma~\ref{lem:countable-mass-tests}, \eqref{eq:integrated-slice} yields
\begin{equation}
 \left|\int_I\sum_{(b,\pi)\in\mathcal T}
                  \partial U_t(b,\pi)\dd t\right|
 \le\int_{\{u\in I\}}|Du|\dd\mu.
 \label{eq:integrated-mass-list}
\end{equation}
Let $\gamma=u_\#(|Du|\mu)$.  Lebesgue differentiation of
\eqref{eq:integrated-mass-list}, simultaneously for the countable
collection of lists on $Y$, gives
\[
 P(t)\le\frac{d\gamma_{\mathrm{ac}}}{dt}(t)
 \qquad\text{for almost every }t.
\]
Here we used \eqref{eq:countable-mass-tests} for each integral slice.
Thus the passage from a signed integrated evaluation to total slice
mass excludes only a countable union of null sets.

Set $\lambda=u_\#\mu$ and $\zeta=u_\#(|Du|^2\mu)$.  For every interval
$I$, Cauchy--Schwarz on $u^{-1}(I)$ and the bound
$|Du|\le\Lip(u)$ imply
\[
 \gamma(I)^2\le\lambda(I)\zeta(I),\qquad
 \gamma(I)\le\Lip(u)\lambda(I).
\]
Divide by the appropriate powers of the interval length and let
symmetric intervals shrink to an almost every point.  The densities
of $\lambda_{\mathrm{ac}}$ and $\zeta_{\mathrm{ac}}$ are respectively
$-a'$ and $-e'$.  Combining the resulting inequalities with the bound
for $P$ proves \eqref{eq:intrinsic-coarea}.
\end{proof}

\subsection{Rearrangement with an arbitrary distribution function}

The small-set perimeter bound and intrinsic coarea now have their
Euclidean coefficients up to the prescribed loss. We construct a radial
Euclidean function with the same distribution of values and controlled
energy. The distribution may have atoms or a singular continuous part;
the generalized inverse below includes both.

\begin{lemma}[Rearranged energy]
\label{lem:rearranged-energy}
Fix $\eta\in(0,1)$ and the corresponding $\delta$ from
Proposition~\ref{prop:small-set}.  If $u\ge0$ is Lipschitz and
$\mu\{u>0\}\le\delta$, there is a nonnegative, radial nonincreasing,
compactly supported Lipschitz function $u^*$ on $\R^n$ such that
\begin{equation}
 \int_{\R^n}(u^*)^q\dd x=\int_Yu^q\dd\mu\quad(q>0),
 \qquad
 \int_{\R^n}|\nabla u^*|^2\dd x
       \le(1-\eta)^{-2}E(u).
 \label{eq:rearranged-energy-bound}
\end{equation}
\end{lemma}

\begin{proof}
If $u=0$ almost everywhere, take $u^*=0$.  Otherwise let
$L=\Lip(u)$ and $U=\esssup_\mu u$.  Then $U>0$ and $L>0$: a function
with Lipschitz constant zero is constant, and a positive constant
would have positivity set of mass $m>\delta$.
For $0\le t\le U$, define
\[
 a(t)=\mu\{u>t\},\qquad
 R(t)=\left(\frac{a(t)}{\omega_n}\right)^{1/n}.
\]
The function $R$ is nonincreasing and right-continuous,
$R(t)>0$ for $t<U$, and $R(U)=0$.
For almost every $t\in(0,U)$, Proposition~\ref{prop:small-set} and
Lemma~\ref{lem:intrinsic-coarea} give
\begin{equation}
 -R'(t)=\frac{-a'(t)}{n\omega_n R(t)^{n-1}}
      \ge c_0,\qquad c_0=\frac{1-\eta}{L}>0.
 \label{eq:radius-distribution-slope}
\end{equation}
This estimate also controls finite differences, even if $R$ has a
singular part.  Indeed, on every compact subinterval of $(0,U)$,
the nonnegative measure $-dR$ has absolutely continuous density
$-R'\ge c_0$ and a nonnegative singular part.  Therefore
\begin{equation}
 R(s)-R(t)\ge c_0(t-s)
 \qquad(0<s<t<U).
 \label{eq:radius-distribution-difference}
\end{equation}

Define the generalized inverse for $\rho\ge0$ by
\begin{equation}
 h(\rho)=\sup\{0\le t<U:\rho<R(t)\},
 \qquad \sup\varnothing:=0.
 \label{eq:generalized-inverse}
\end{equation}
It is nonincreasing, $h(0)=U$, and $h(\rho)=0$ for
$\rho\ge R(0)$.  It is also $c_0^{-1}$-Lipschitz.  To check the latter
claim, let $\rho_1<\rho_2$ with $h(\rho_1)>h(\rho_2)$ and choose
\[
 h(\rho_2)<s<t<h(\rho_1).
\]
The definition and monotonicity imply
$R(s)\le\rho_2$ and $R(t)>\rho_1$, whence
$c_0(t-s)\le R(s)-R(t)\le\rho_2-\rho_1$.
Letting $s$ and $t$ approach the two inverse values proves the claim.

For $0<t<U$, right continuity of $R$ implies
\begin{equation}
 h(\rho)>t\quad\Longleftrightarrow\quad \rho<R(t).
 \label{eq:inverse-superlevels}
\end{equation}
In the reverse implication, if $\rho<R(t)$, right continuity supplies
some $s>t$ with $\rho<R(s)$; the other implication follows from
monotonicity.  Consequently $u^*(x)=h(|x|)$ is Lipschitz and compactly
supported, and its superlevel sets have Euclidean volume $a(t)$.
The layer-cake formula proves the asserted equality of all positive
power integrals.

We give the exact energy substitution, including possible atoms and
singular continuous parts of $a$.  Strict decrease in
\eqref{eq:radius-distribution-difference} implies
\begin{equation}
 h(R(t))=t\qquad(0<t<U).
 \label{eq:inverse-identity}
\end{equation}
Indeed, all levels $s<t$ satisfy $R(s)>R(t)$, while no $s\ge t$ is
eligible in \eqref{eq:generalized-inverse} with $\rho=R(t)$.
If $R$ is continuous at $t$, \eqref{eq:inverse-superlevels} and this
continuity show that the entire fiber $h^{-1}(\{t\})$ is the singleton
$\{R(t)\}$.

The monotone function $R$ is continuous except at countably many
points and differentiable with finite derivative almost everywhere.
Its derivative, wherever it is relevant to
\eqref{eq:radius-distribution-slope}, is nonzero.
Let $N$ be the null set where the Lipschitz function $h$ is not
differentiable.  The set $h(N)$ is null because Lipschitz functions
map Lebesgue null sets to null sets.  By \eqref{eq:inverse-identity},
every level $t$ with $R(t)\in N$ lies in $h(N)$.  It follows that for
almost every $t\in(0,U)$ both derivatives needed below exist, and
differentiating the inverse identity gives
\begin{equation}
 h'(R(t))=\frac1{R'(t)}.
 \label{eq:inverse-derivative}
\end{equation}
For example, this follows directly by letting $s\to t$ in
$h(R(s))-h(R(t))=s-t$ at a continuity and differentiability point
of $R$.

Since $h$ is absolutely continuous and nonincreasing on $[0,R(0)]$,
the measure $|h'(\rho)|\dd\rho$ pushes forward under $h$ to Lebesgue
measure on $(0,U)$.  One may verify this first on an interval
$(a,b)\subset(0,U)$: integration of $-h'$ over its inverse image
equals $b-a$ by the fundamental theorem of calculus; constant
portions have zero derivative.  Intervals determine the measure.
At almost every target level the fiber consists of the single point
$R(t)$, so this substitution and \eqref{eq:inverse-derivative} give
\begin{align}
 \int_{\R^n}|\nabla u^*|^2\dd x
 &=n\omega_n\int_0^{R(0)}\rho^{n-1}|h'(\rho)|^2\dd\rho \notag\\
 &=\int_0^U\frac{n\omega_n R(t)^{n-1}}{-R'(t)}\dd t \notag\\
 &=\int_0^U
       \frac{\bigl(n\omega_n R(t)^{n-1}\bigr)^2}{-a'(t)}\dd t.
 \label{eq:rearranged-energy-identity}
\end{align}
The denominators are positive and finite almost everywhere by
\eqref{eq:radius-distribution-slope}.  The preceding substitution
also justifies the equality for a nonnegative integrand before its
finiteness is known.

For clarity, a jump of $R$ produces an interval on which $h$ is
constant.  In particular, an atom at $U$ produces a central constant
part of $u^*$.  Such intervals contribute no energy.  A singular
continuous decrease causes no additional term either: the exact
substitution uses the absolutely continuous measure
$|h'|\dd\rho$, and the inverse image of any null set of levels has
zero measure for this measure.  Thus no absolute-continuity
assumption on $a$ has entered
\eqref{eq:rearranged-energy-identity}.

Finally, \eqref{eq:small-set} gives
\[
 \bigl(n\omega_n R(t)^{n-1}\bigr)^2
       \le(1-\eta)^{-2}P(t)^2.
\]
Using \eqref{eq:intrinsic-coarea} in
\eqref{eq:rearranged-energy-identity} therefore yields
\[
 \int_{\R^n}|\nabla u^*|^2\dd x
   \le(1-\eta)^{-2}\int_0^U(-e'(t))\dd t
   \le(1-\eta)^{-2}E(u).
\]
The last inequality only uses monotonicity of $e$ and remains true
if its distributional derivative has a singular part.
\end{proof}

\subsection{The sharp radial Euclidean inequality}

Rearrangement reduces the small-support estimate to a radial function
on $\R^n$. The next lemma identifies its sharp energy coefficient;
afterwards only the truncation to arbitrary functions remains.
This is the radial case of the sharp Sobolev inequality of Aubin and
Talenti~\cite{Aubin1976,Talenti1976}. Its proof is an explicit radial
version of the mass-transport argument of Cordero-Erausquin, Nazaret and
Villani~\cite[Section~2]{CorderoNazaretVillani2004}: cumulative masses
determine the transport, and determinant comparison followed by
integration by parts recovers the sharp coefficient.

\begin{lemma}[Sharp radial Sobolev inequality]
\label{lem:radial-sobolev}
For every nonnegative, radial nonincreasing, compactly supported
Lipschitz function $f$ on $\R^n$, where $n>2$,
\begin{equation}
 ns_n^{2/n}
       \left(\int_{\R^n}f^p\dd x\right)^{2/p}
    \le \frac4{n-2}\int_{\R^n}|\nabla f|^2\dd x.
 \label{eq:radial-sobolev}
\end{equation}
\end{lemma}

\begin{proof}
The assertion is immediate for $f=0$.  By homogeneity, assume
$\int f^p\dd x=1$.  Write $f=f(r)$, $r=|x|$.  Continuity and radial
monotonicity give a finite $r_0>0$ such that $f>0$ on $[0,r_0)$ and
$f=0$ on $[r_0,\infty)$.
For $D>0$ define
\[
 g_D(r)=A_D(1+r^2)^{-(n-2)/2}\quad(0\le r<D),
 \qquad g_D(r)=0\quad(r\ge D),
\]
where $A_D>0$ is chosen so that $\int g_D^p\dd x=1$.
For $0<r<r_0$ match the cumulative radial masses by
\begin{equation}
 n\omega_n\int_0^r s^{n-1}f(s)^p\dd s
   =n\omega_n\int_0^{t(r)}s^{n-1}g_D(s)^p\dd s.
 \label{eq:radial-transport-cdf}
\end{equation}
Both cumulative functions are $C^1$ with strictly positive derivative
on their open radial intervals.  Their inverse composition $t$ is
therefore $C^1$ on $(0,r_0)$, strictly increasing, and satisfies
$t(0+)=0$ and $t(r_0-)=D$.  Differentiation gives
\begin{equation}
 d_T:=t'(r)\left(\frac{t(r)}r\right)^{n-1}
      =\frac{f(r)^p}{g_D(t(r))^p}.
 \label{eq:radial-transport-jacobian}
\end{equation}
The radial map $T(x)=t(|x|)x/|x|$ has positive eigenvalues
$t'$ and $n-1$ copies of $t/r$ and transports $f^p\dd x$ to
$g_D^p\dd x$.  This follows either from
\eqref{eq:radial-transport-cdf} by radial integration, or from the
change of variables and \eqref{eq:radial-transport-jacobian}.
The jump of $g_D$ at $D$ is harmless: the calculation uses its
positive continuous interior profile and never differentiates the
cutoff.

Let
\[
 \ell=p\left(1-\frac1n\right)=\frac{2(n-1)}{n-2},
 \qquad \ell-1=\frac p2.
\]
Change of variables, the arithmetic--geometric mean inequality for
the positive eigenvalues of $DT$, and radial integration by parts
give
\begin{align}
 n\int_{\R^n}g_D^\ell\dd x
 &=n\int_{|x|<r_0}f^\ell d_T^{1/n}\dd x \notag\\
 &\le\int_{|x|<r_0}f^\ell
               \left(t'+(n-1)\frac tr\right)\dd x \notag\\
 &=-\ell\int_{|x|<r_0}f^{p/2}f't\dd x \notag\\
 &\le\ell\left(\int_{\R^n}|\nabla f|^2\dd x\right)^{1/2}
          \left(\int_{\R^n}|x|^2g_D^p\dd x\right)^{1/2}.
 \label{eq:radial-transport-estimate}
\end{align}
The last equality needed for Cauchy--Schwarz is the transport
identity $\int f^pt^2\dd x=\int|x|^2g_D^p\dd x$.
To justify the integration by parts, first integrate over
$\varepsilon<r<R<r_0$.  The boundary term is
$n\omega_n[r^{n-1}tf^\ell]_{\varepsilon}^R$.
It tends to zero at $r_0$ because $f(r_0)=0$ and $t\le D$, and at zero
because $n>1$ and $t,f$ are bounded.  The resulting derivative
integral is absolutely integrable: $f'$ is essentially bounded,
$t$ is bounded, and the source ball is bounded.  The source profile
is Lipschitz, so its radial integration by parts is justified by
absolute continuity.

Now put
\[
 I_0=\int_{\R^n}(1+|x|^2)^{-n}\dd x,
 \qquad A_\infty=I_0^{-1/p},\qquad
 g(r)=A_\infty(1+r^2)^{-(n-2)/2}.
\]
The normalizing constants $A_D$ converge to $A_\infty$ as
$D\to\infty$.  Moreover,
\[
 \int g_D^\ell\dd x\longrightarrow\int g^\ell\dd x,
 \qquad
 \int|x|^2g_D^p\dd x\longrightarrow\int|x|^2g^p\dd x.
\]
For both integrals the unnormalized radial integrand is of order
$r^{1-n}$ at infinity.  Thus both are finite for every $n>2$, and
the convergence follows by monotone convergence for the
unnormalized integrals and convergence of $A_D$.
In particular, no limit of the transport maps is required.

The ratio on the limiting left side of
\eqref{eq:radial-transport-estimate} can be evaluated from $g$ itself.
Direct differentiation shows that
\[
 -g'(r)=Krg(r)^{p/2},\qquad
 K=(n-2)A_\infty^{-2/(n-2)}>0.
\]
If $J=\int g^\ell\dd x$ and $M_2=\int|x|^2g^p\dd x$, radial
integration by parts with the vector field $x$ gives
\[
 nJ=-\ell\int g^{p/2}g'r\dd x=\ell K M_2,
 \qquad \int|\nabla g|^2\dd x=K^2M_2.
\]
Its boundary term $r^ng(r)^\ell$ vanishes at zero and is
$O(r^{2-n})$ at infinity.  Therefore
\begin{equation}
 \frac{nJ}{\ell M_2^{1/2}}
    =\left(\int|\nabla g|^2\dd x\right)^{1/2}.
 \label{eq:bubble-transport-ratio}
\end{equation}
Taking $D\to\infty$ in \eqref{eq:radial-transport-estimate} proves
$\int|\nabla f|^2\dd x\ge\int|\nabla g|^2\dd x$.

The transport argument has reduced the lower energy bound to the
single comparison profile $g$. It remains to compute this energy and
thus fix the threshold used in the current's critical minimization.
The inverse stereographic map
\[
 x\longmapsto\frac{(2x,|x|^2-1)}{1+|x|^2}\in\mathbb S^n
\]
pulls the sphere metric back to $4(1+|x|^2)^{-2}$ times the Euclidean
metric.  Its volume Jacobian is $2^n(1+|x|^2)^{-n}$, so
\begin{equation}
 I_0=2^{-n}s_n.
 \label{eq:stereographic-integral}
\end{equation}
For $w(r)=(1+r^2)^{-(n-2)/2}$, direct differentiation gives
\[
 -\Delta w=n(n-2)(1+r^2)^{-(n+2)/2}.
\]
Integration by parts consequently yields
\begin{align*}
 \int_{\R^n}|\nabla g|^2\dd x
 &=n(n-2)A_\infty^2I_0\\
 &=n(n-2)I_0^{1-2/p}
   =\frac{n(n-2)}4s_n^{2/n}.
\end{align*}
Here the boundary term $r^{n-1}g(r)g'(r)$ vanishes at zero and is
$O(r^{2-n})$ at infinity; also $1-2/p=2/n$.
Multiplying the energy lower bound by $4/(n-2)$ proves
\eqref{eq:radial-sobolev} for the normalized $f$, and homogeneity
proves the general statement.  All moments and boundary limits
used above remain valid when $n=3$ or $n=4$; no $L^2$ integrability
of the full profile $g$ is needed.

The coefficient is also optimal within the class in the statement.
Indeed, choose a nonnegative nonincreasing Lipschitz cutoff $\chi$
equal to one on $[0,1]$ and zero on $[2,\infty)$, and set
$g_R(x)=g(x)\chi(|x|/R)$.  Then $g_R\to g$ in $L^p$, and
$\nabla g_R\to\nabla g$ in $L^2$: the tail of $\nabla g$ tends to
zero, while the extra cutoff energy is bounded by
\[
 \frac{\Lip(\chi)^2}{R^2}
       \int_{R<|x|<2R}g(x)^2\dd x=O(R^{2-n})\longrightarrow0.
\]
Each $g_R$ is admissible, and its Sobolev quotient tends to the
computed value for $g$.
\end{proof}

\subsection{Globalization}

The preceding two lemmas give the sharp coefficient, with an
arbitrarily small loss, for functions supported on a small amount of
current mass. Truncation at a level controlled by $\|u\|_2$ yields the
section's asserted inequality for every Lipschitz function.

\begin{proposition}[Almost Euclidean critical Sobolev bound]
\label{prop:almost-sobolev}
For every $\epsilon\in(0,S_*)$ there is a finite constant
$C_\epsilon\ge0$ such that every real Lipschitz function $u$ on $Y$
satisfies
\begin{equation}
 (S_*-\epsilon)\|u\|_p^2
     \le4\beta E(u)+C_\epsilon\|u\|_2^2.
 \label{eq:almost-sobolev}
\end{equation}
The constant is independent of $u$.
\end{proposition}

\begin{proof}
Fix temporarily $\eta\in(0,1)$ and its corresponding $\delta$.
Combining Lemmas~\ref{lem:rearranged-energy} and
\ref{lem:radial-sobolev} gives
\begin{equation}
 S_*(1-\eta)^2\|z\|_p^2\le4\beta E(z)
 \quad\text{if }z\ge0\text{ is Lipschitz and }\mu\{z>0\}\le\delta.
 \label{eq:small-support-sobolev}
\end{equation}
This also holds when $z=0$ almost everywhere.

For a general Lipschitz $u$ with $\|u\|_2>0$, set
\[
 \tau=\frac{\|u\|_2}{\sqrt\delta},\qquad z=(|u|-\tau)_+.
\]
Chebyshev's inequality gives $\mu\{z>0\}\le\delta$.
Scalar Lipschitz composition in the charts gives $E(z)\le E(u)$;
at the finitely many corner levels the gradient vanishes almost
everywhere on the corresponding level sets, so the usual
almost-everywhere chain rule applies.  In addition,
$0\le|u|-z\le\tau$, whence
\[
 \||u|-z\|_p\le m^{1/p}\tau.
\]
For every $\xi>0$, the triangle inequality followed by
$(b+d)^2\le(1+\xi)b^2+(1+\xi^{-1})d^2$ now yields
\begin{align*}
 \frac{S_*(1-\eta)^2}{1+\xi}\|u\|_p^2
 &\le S_*(1-\eta)^2\|z\|_p^2
       +\frac{S_*(1-\eta)^2}{\xi}m^{2/p}\tau^2\\
 &\le4\beta E(u)
       +\frac{S_*(1-\eta)^2m^{2/p}}{\xi\delta}\|u\|_2^2.
\end{align*}
Choose $\eta>0$ and $\xi>0$ sufficiently small that
$S_*(1-\eta)^2/(1+\xi)\ge S_*-\epsilon$, and then fix the associated
$\delta$.  The last display proves the assertion with, for example,
\[
 C_\epsilon=\frac{S_*(1-\eta)^2m^{2/p}}{\xi\delta}.
\]
If $\|u\|_2=0$, then $\|u\|_p=0$ and the asserted inequality is
immediate.
\end{proof}
